%------------------------------------------------------------------------------%
\section{Relações no Triângulo Retângulo}

Importantes relações são estabelecidas entre os ângulos e os lados de
triângulos retângulos. Considere um triângulo retângulo que contém um ângulo
$0<\theta<\pi/2$ conforme Figura~\ref{fig:triangulo-retangulo}.

%------------------------------------------------------------------------------%
\begin{figure}
  \centering
  \includegraphics[width=0.5\textwidth,trim={1cm 20cm 19cm 9cm},clip]{figdeflados}
  \caption{Triangulo retângulo}
  \label{fig:triangulo-retangulo}
\end{figure}
%------------------------------------------------------------------------------%

Em relação ao ângulo $\theta$, os lados desse triângulo
recebem os seguintes nomes:
\begin{itemize}
  \item \emph{cateto oposto}\index{Cateto!oposto}
  \item \emph{cateto }\index{Cateto!adjacente}
  \item \emph{hipotenusa}\index{Hipotenusa}
\end{itemize}

Como a soma dos ângulos internos de qualquer triângulo é $\pi$ radianos, o
terceiro ângulo (que chamaremos de $\alpha$) desse triângulo mede
$\pi/2-\theta$ (pois $\alpha+\theta+\pi/2=\pi$). Portanto, o cateto oposto ao
ângulo $\theta$ é o cateto adjacente ao ângulo $\alpha$ e o cateto adjacente ao
ângulo $\theta$ é o cateto oposto ao ângulo $\alpha$.

Estudaremos as relações entre os lados de dois triângulos retângulos que
possuem um mesmo ângulo $\theta$. Fixamos dois desses triângulos e nomeamos os
lados desses triângulos da seguinte forma:

%------------------------------------------------------------------------------%
\begin{figure}
  \centering
  \includegraphics[width=0.5 \textwidth,trim={1cm 28cm 17cm 1cm} ,clip]{figdeffunc}
  \caption{?}
\end{figure}
%------------------------------------------------------------------------------%

Como ambos os triângulos possuem um ângulo reto e um ângulo $\theta$, o
terceiro ângulo desses dois triângulos mede $\pi/2-\theta$. Ou seja, eles são
semelhantes pelo caso $AAA$ (ângulo ângulo ângulo) e podemos estabelecer as
seguintes relações:
\[
  \frac{y_1}{z_1} = \frac{y_2}{z_2},  \qquad
  \frac{x_1}{z_1} = \frac{x_2}{z_2},  \qquad
  \frac{y_1}{x_1} = \frac{y_2}{x_2}.
\]
Ou seja, independente dos valores dos catetos e da hipotenusa, as razões entre
os lados de um triângulo retângulo que contém um ângulos $\theta$ são
constantes. Essas constantes (que dependem de $\theta$) recebem nomes, os
principais deles são: (Outras razões recebem os nomes de secante,
cosseno e cotangente.)

\emph{Seno de $\theta$}
\[
  \sin\theta=\frac{\text{cateto oposto}}{\text{hipotenusa}}.
\]

\emph{Cosseno de $\theta$}
\[
  \cos\theta=\frac{\text{cateto adjacente}}{\text{hipotenusa}}.
\]

\emph{Tangente de $\theta$}
\[
  \tan\theta=\frac{\text{cateto oposto}}{\text{cateto adjacente}}.
\]

Assim definidos, seno, cosseno e tangente são funções cujo domínio é o
intervalo \(\left(0, \sfrac{\pi}{2}\right)\).
Para estendermos essas definições e obtermos funções com
domínio real precisaremos generalizar a noção de ângulos para valores maiores
que $2\pi$ e ângulos negativos e posteriormente, estender a noção de seno,
cosseno e tangente desses ângulos através do círculo trigonométrico unitário.

%------------------------------------------------------------------------------%
\begin{observation}
  É possível calcular os valores de seno, cosseno e tangente para os ângulos
  \(\sfrac{\pi}{4}\),
  \(\sfrac{\pi}{3}\) e
  \(\sfrac{\pi}{6}\)
  utilizando a relação de Pitágoras no quadrado de
  lado 1 e no triângulo equilátero de lado 1. Esses ângulos são chamados de
  notáveis porque para demais ângulos não existe uma maneira tão simples de
  estabelecer os valores de seno, cosseno e tangente.
\end{observation}
%------------------------------------------------------------------------------%

%------------------------------------------------------------------------------%
\begin{example}
  \label{ex:angulos_notaveis}
  Calcule os valores de seno, cosseno e tangente para os ângulos
  \(\sfrac{\pi}{4}\),
  \(\sfrac{\pi}{3}\) e
  \(\sfrac{\pi}{6}\)
  utilizando a relação de Pitágoras no quadrado de lado 1 e
  no triângulo equilátero de lado 1, justificando a tabela a seguir.

  \solution
  Vamos mostrar as razões trigonométricas para o ângulo de \ang{45}. Para
  isso, considere um quadrado de lado 1 e tracemos sua diagonal. Ao fazermos
  esse procedimento, dividimos o quadrado em dois triângulos retângulos iguais
  como mostra a figura.

  \IncludeGraphicCentredEx{quadrado}[0.4][-1][-1]

  Podemos observar que
  \begin{description}
      \item[] \qquad Cateto oposto a    \ang{45} vale 1;
      \item[] \qquad Cateto adjacente a \ang{45} vale 1;
      \item[] \qquad Hipotenusa desses triângulos retângulos é $\sqrt{2}$.
  \end{description}
  Assim,
  \begin{align*}
     \sin(\ang{45}) & = \frac{1}{\sqrt{2}} = \frac{\sqrt{2}}{2}; \\[2mm]
     \cos(\ang{45}) & = \frac{1}{\sqrt{2}} = \frac{\sqrt{2}}{2}; \\[2mm]
     \tan(\ang{45}) & = \frac{1}{1}        = 1.
  \end{align*}
  Para encontrarmos as razões trigonométricas para os ângulos de \ang{30} e
  \ang{60} basta traçarmos uma altura (que é mediana e bissetriz) em um
  triângulo equilátero de lado $l$.
\end{example}
%------------------------------------------------------------------------------%

%------------------------------------------------------------------------------%

A partir dos valores encontrados no Exemplo~\ref{ex:angulos_notaveis},
construímos a seguinte tabela
\begin{ceqn}
\[
  \begin{array}{c<{\quad}>{\quad}c<{\quad}>{\quad}c<{\quad}>{\quad}c}  \hline
    \theta          & \sin\theta          & \cos\theta          & \tan\theta          \\ \hline
    \dfrac{\pi}{6}  & \dfrac{1}{2}        & \dfrac{\sqrt{3}}2{} & \dfrac{\sqrt{3}}{3} \\[3mm]
    \dfrac{\pi}{4}  & \dfrac{\sqrt{2}}{2} & \dfrac{\sqrt{2}}{2} &      1              \\[3mm]
    \dfrac{\pi}{3}  & \dfrac{\sqrt{3}}{2} & \dfrac{1}{2}        &  \sqrt{3}           \\[3mm] \hline
  \end{array}
\]
\end{ceqn}

A seguir resolveremos alguns exercícios que envolvem as razões métricas em um
triângulo retângulo.

%------------------------------------------------------------------------------%
\begin{example} \label{trigon1ex}
  Um avião decola, percorrendo uma trajetória retilínea, formando com o solo,
  um ângulo de \ang{30} (suponha que a região sobrevoada pelo avião seja
  plana). Após percorrer 1000 metros, qual a altura atingida pelo avião?

  \solution
  A figura descreve a situação.

  \IncludeGraphicCentredEx{exerc1t}[0.5]

  Note que temos envolvidos no problema um ângulo (\ang{30}), o cateto
  oposto a esse ângulo ($h$) e a hipotenusa (\SI{1000}{\meter}), assim,
  é conveniente usar o seno para obtermos a medida desconhecida.
  \begin{align*}
    \sin(\ang{30}) & = \frac{h}{1000}   \\
    \frac{1}{2}    & = \frac{h}{1000}   \\
    h              & = \SI{500}{\meter}
  \end{align*}
  Assim, após percorrer \SI{1000}{\meter} o avião atinge \SI{500}{\meter}.
\end{example}
%------------------------------------------------------------------------------%

%------------------------------------------------------------------------------%
\begin{figure}
  \centering
  \caption{Figura para o exercício \ref{exerc34}.}
\end{figure}
%------------------------------------------------------------------------------%

%------------------------------------------------------------------------------%
\begin{example} \label{exerc34}
  Um sensor detecta a aproximação de veículos em uma rua. Determine uma função
  $x(\theta)$ que forneça a distância do sensor ao veículo em relação a
  $\theta$. Se o sensor só detecta objetos quando $\theta$ é menor ou igual a
  $\sfrac{\pi}{3}$,
  determine a distância máxima que o carro pode estar do sensor para
  ser detectado.

  \solution
  \IncludeGraphicCentredEx{figexercar}[0.5]

  A distância máxima ocorre quando o ângulo no qual o sensor detecta os objetos
  é máxima. Como, no exercício, temos o valor do cateto adjacente ao ângulo
  $\theta$ e queremos encontrar a hipotenusa. Portanto, iremos usar o cosseno.
  \begin{align*}
    \cos(\ang{60}) & = \frac{8}{x}     \\[2mm]
    \frac{1}{2}    & = \frac{8}{x}     \\[2mm]
    x              & = \SI{16}{\meter}
  \end{align*}
  Portanto, a distância máxima que o carro pode estar do sensor para ser
  detectado é de \SI{16}{\meter}.
\end{example}
%------------------------------------------------------------------------------%
